\documentclass[10pt,a4paper]{amsart}
\usepackage{amsmath,amssymb,amsthm,graphicx}
% For commutative diagrams we use Xy-pic
\usepackage[all]{xy}

% We use 2cell for 2-commutative diagrams.
\xyoption{2cell}
% Theorem environments.
%
\theoremstyle{plain}
\newtheorem{theorem}[subsection]{Theorem}
\newtheorem{proposition}[subsection]{Proposition}
\newtheorem{lemma}[subsection]{Lemma}

\theoremstyle{definition}
\newtheorem{definition}[subsection]{Definition}
\newtheorem{example}[subsection]{Example}
\newtheorem{exercise}[subsection]{Exercise}
\newtheorem{situation}[subsection]{Situation}

\theoremstyle{remark}
\newtheorem{remark}[subsection]{Remark}
\newtheorem{remarks}[subsection]{Remarks}

\numberwithin{equation}{subsection}

% Macros
%
\def\lim{\mathop{\mathrm{lim}}\nolimits}
\def\colim{\mathop{\mathrm{colim}}\nolimits}
\def\Spec{\mathop{\mathrm{Spec}}}
\def\Hom{\mathop{\mathrm{Hom}}\nolimits}
\def\Ext{\mathop{\mathrm{Ext}}\nolimits}
\def\SheafHom{\mathop{\mathcal{H}\!\mathit{om}}\nolimits}
\def\SheafExt{\mathop{\mathcal{E}\!\mathit{xt}}\nolimits}
\def\Sch{\mathit{Sch}}
\def\Mor{\mathop{\mathrm{Mor}}\nolimits}
\def\Ob{\mathop{\mathrm{Ob}}\nolimits}
\def\Sh{\mathop{\mathit{Sh}}\nolimits}
\def\NL{\mathop{N\!L}\nolimits}
\def\CH{\mathop{\mathrm{CH}}\nolimits}
\def\proetale{{pro\text{-}\acute{e}tale}}
\def\etale{{\acute{e}tale}}
\def\QCoh{\mathit{QCoh}}
\def\Ker{\mathop{\mathrm{Ker}}}
\def\Im{\mathop{\mathrm{Im}}}
\def\Coker{\mathop{\mathrm{Coker}}}
\def\Coim{\mathop{\mathrm{Coim}}}

% Boxtimes
%
\DeclareMathSymbol{\boxtimes}{\mathbin}{AMSa}{"02}

%
% Macros for moduli stacks/spaces
%
\def\QCohstack{\mathcal{QC}\!\mathit{oh}}
\def\Cohstack{\mathcal{C}\!\mathit{oh}}
\def\Spacesstack{\mathcal{S}\!\mathit{paces}}
\def\Quotfunctor{\mathrm{Quot}}
\def\Hilbfunctor{\mathrm{Hilb}}
\def\Curvesstack{\mathcal{C}\!\mathit{urves}}
\def\Polarizedstack{\mathcal{P}\!\mathit{olarized}}
\def\Complexesstack{\mathcal{C}\!\mathit{omplexes}}
% \Pic is the operator that assigns to X its picard group, usage \Pic(X)
% \Picardstack_{X/B} denotes the Picard stack of X over B
% \Picardfunctor_{X/B} denotes the Picard functor of X over B
\def\Pic{\mathop{\mathrm{Pic}}\nolimits}
\def\Picardstack{\mathcal{P}\!\mathit{ic}}
\def\Picardfunctor{\mathrm{Pic}}
\def\Deformationcategory{\mathcal{D}\!\mathit{ef}}
\usepackage{cleveref}
\crefname{originallemma}{original lemma}{original lemmas}
\Crefname{originallemma}{Original lemma}{Original lemmas}
\makeatletter
\newlabel{divisors-definition-pullback-effective-Cartier-divisor}{{01WV}{source}{}{sourcepointer-1}{}}
\newlabel{divisors-definition-pullback-effective-Cartier-divisor@cref}{{[originallemma][0][]01WV}{[1][0][]source}{}{}{}}
\makeatother
\crefname{originallemma}{original lemma}{original lemmas}
\Crefname{originallemma}{Original lemma}{Original lemmas}
\makeatletter
\newlabel{divisors-lemma-pullback-effective-Cartier-defined}{{02OO}{source}{}{sourcepointer-2}{}}
\newlabel{divisors-lemma-pullback-effective-Cartier-defined@cref}{{[originallemma][0][]02OO}{[1][0][]source}{}{}{}}
\makeatother
\crefname{originallemma}{original lemma}{original lemmas}
\Crefname{originallemma}{Original lemma}{Original lemmas}
\makeatletter
\newlabel{lemma-cycle-push-sheaf}{{02R6}{source}{}{sourcepointer-3}{}}
\newlabel{lemma-cycle-push-sheaf@cref}{{[originallemma][0][]02R6}{[1][0][]source}{}{}{}}
\makeatother
\crefname{originallemma}{original lemma}{original lemmas}
\Crefname{originallemma}{Original lemma}{Original lemmas}
\makeatletter
\newlabel{lemma-divisor-delta-dimension}{{02RK}{source}{}{sourcepointer-4}{}}
\newlabel{lemma-divisor-delta-dimension@cref}{{[originallemma][0][]02RK}{[1][0][]source}{}{}{}}
\makeatother
\crefname{originallemma}{original lemma}{original lemmas}
\Crefname{originallemma}{Original lemma}{Original lemmas}
\makeatletter
\newlabel{lemma-proper-pushforward-alteration}{{02RT}{source}{}{sourcepointer-5}{}}
\newlabel{lemma-proper-pushforward-alteration@cref}{{[originallemma][0][]02RT}{[1][0][]source}{}{}{}}
\makeatother
\crefname{originallemma}{original lemma}{original lemmas}
\Crefname{originallemma}{Original lemma}{Original lemmas}
\makeatletter
\newlabel{morphisms-lemma-composition-proper}{{01W3}{source}{}{sourcepointer-6}{}}
\newlabel{morphisms-lemma-composition-proper@cref}{{[originallemma][0][]01W3}{[1][0][]source}{}{}{}}
\makeatother
\crefname{originallemma}{original lemma}{original lemmas}
\Crefname{originallemma}{Original lemma}{Original lemmas}
\makeatletter
\newlabel{situation-setup}{{02QL}{source}{}{sourcepointer-7}{}}
\newlabel{situation-setup@cref}{{[originallemma][0][]02QL}{[1][0][]source}{}{}{}}
\makeatother
\title{The closure of the graph of a rational function in a projective line bundle: properness, the two pullback divisors and the divisor of f}
\author{The Stacks Project Authors}
\date{stacks-project:chow.tex:local-snapshot-20261012; GNU FDL 1.2 or later}
\begin{document}
\maketitle
\noindent\textit{Editorial note.} All delivered mathematical text is the author's own, copied from the pinned Stacks Project chapter file chow.tex; no statement or proof step was rewritten, reordered or removed.\par\medskip
\noindent\textit{Editorial note.} Preamble provenance: the pinned chapter file opens with \textbackslash{}input\{preamble\} and carries no preamble text of its own, so the project's single Stacks-wide preamble staged at sources/stacks-chow/preamble.tex is excerpted verbatim as two spans, lines 14-18 (Xy-pic setup) and lines 154-217 (the native theorem declarations and mathematical macros). Nothing was added, removed or redefined; the wrapper class and the amsmath/amssymb/amsthm/graphicx base packages are supplied by the assembler.\par\medskip
\noindent\textit{Editorial note.} Same-chapter pointers are located at their \textbackslash{}label lines inside the pinned chapter file; cross-chapter pointers are located in the project's own tag list staged in the snapshot (sources/stacks-chow/tags.txt), at the verbatim lines carrying their canonical tags.\par\medskip
\noindent\textit{Allowed standard prerequisites (editorial).}\begin{itemize}
\item Composition of proper morphisms. Morphisms, label morphisms-lemma-composition-proper, invoked at proof line 2746. Makes the projection p : Y -> X proper as the composite of the closed immersion Y -> P\textasciicircum{}1\_X with the proper projection P\textasciicircum{}1\_X -> X; this is assertion (1).
\item Pullback of an effective Cartier divisor along a morphism. Divisors, labels divisors-definition-pullback-effective-Cartier-divisor (statement, proof line 2716) and divisors-lemma-pullback-effective-Cartier-defined (proof line 2756). Defines Y\_0 and Y\_infty as pullbacks of the two standard sections and supplies the local criterion that makes them defined because q(eta') avoids D\_0 and D\_infty; assertions (3) and the dimension count following it are built on this.
\item Dimension of the support of a divisor and compatibility of pushforward with divisors of rational functions. Same chapter chow.tex, labels lemma-divisor-delta-dimension (proof line 2759), lemma-cycle-push-sheaf (proof line 2790) and lemma-proper-pushforward-alteration (proof line 2794). Gives the delta-dimension of Y\_0 and Y\_infty, identifies p\_*[Y\_0]\_\{n-1\} with the class of the image of Y\_0, and supplies the compatibility of p\_* with div(f) that yields assertion (5).
\end{itemize}\medskip
\ifdefined\hypertarget\hypertarget{sourcepointer-1}{}\fi
\noindent\textit{Source reference.} Original lemma 01WV: sources/stacks-chow/tags.txt, source lines 2496--2496. Navigation pointer for \textbackslash{}ref\{divisors-definition-pullback-effective-Cartier-divisor\} used inside the counted span (line 2716); the argument that consumes it is retained in the span.\par\medskip
\ifdefined\hypertarget\hypertarget{sourcepointer-2}{}\fi
\noindent\textit{Source reference.} Original lemma 02OO: sources/stacks-chow/tags.txt, source lines 3523--3523. Navigation pointer for \textbackslash{}ref\{divisors-lemma-pullback-effective-Cartier-defined\} used inside the counted span (line 2756); the argument that consumes it is retained in the span.\par\medskip
\ifdefined\hypertarget\hypertarget{sourcepointer-3}{}\fi
\noindent\textit{Source reference.} Original lemma 02R6: sources/stacks-chow/pinned-original-source.tex, source lines 1992--1992. Navigation pointer for \textbackslash{}ref\{lemma-cycle-push-sheaf\} used inside the counted span (line 2790); the argument that consumes it is retained in the span.\par\medskip
\ifdefined\hypertarget\hypertarget{sourcepointer-4}{}\fi
\noindent\textit{Source reference.} Original lemma 02RK: sources/stacks-chow/pinned-original-source.tex, source lines 2489--2489. Navigation pointer for \textbackslash{}ref\{lemma-divisor-delta-dimension\} used inside the counted span (line 2759); the argument that consumes it is retained in the span.\par\medskip
\ifdefined\hypertarget\hypertarget{sourcepointer-5}{}\fi
\noindent\textit{Source reference.} Original lemma 02RT: sources/stacks-chow/pinned-original-source.tex, source lines 2645--2645. Navigation pointer for \textbackslash{}ref\{lemma-proper-pushforward-alteration\} used inside the counted span (line 2794); the argument that consumes it is retained in the span.\par\medskip
\ifdefined\hypertarget\hypertarget{sourcepointer-6}{}\fi
\noindent\textit{Source reference.} Original lemma 01W3: sources/stacks-chow/tags.txt, source lines 2468--2468. Navigation pointer for \textbackslash{}ref\{morphisms-lemma-composition-proper\} used inside the counted span (line 2746); the argument that consumes it is retained in the span.\par\medskip
\ifdefined\hypertarget\hypertarget{sourcepointer-7}{}\fi
\noindent\textit{Source reference.} Original lemma 02QL: sources/stacks-chow/pinned-original-source.tex, source lines 1372--1372. Navigation pointer for \textbackslash{}ref\{situation-setup\} used inside the counted span (line 2702); the argument that consumes it is retained in the span.\par\medskip
\begin{lemma}
\label{lemma-rational-function}
Let $(S, \delta)$ be as in Situation \ref{situation-setup}.
Let $X$ be locally of finite type over $S$. Assume $X$ is
integral and $n = \dim_\delta(X)$. Let $f \in R(X)^*$.
Let $U \subset X$ be a nonempty open such that $f$
corresponds to a section $f \in \Gamma(U, \mathcal{O}_X^*)$.
Let $Y \subset X \times_S \mathbf{P}^1_S$ be the
closure of the graph of $f : U \to \mathbf{P}^1_S$.
Then
\begin{enumerate}
\item the projection morphism $p : Y \to X$ is proper,
\item $p|_{p^{-1}(U)} : p^{-1}(U) \to U$ is an isomorphism,
\item the pullbacks $Y_0 = q^{-1}D_0$ and $Y_\infty = q^{-1}D_\infty$
via the morphism $q : Y \to \mathbf{P}^1_S$ are defined
(Divisors, Definition
\ref{divisors-definition-pullback-effective-Cartier-divisor}),
\item we have
$$
\text{div}_Y(f) = [Y_0]_{n - 1} - [Y_\infty]_{n - 1}
$$
\item we have
$$
\text{div}_X(f) = p_*\text{div}_Y(f)
$$
\item if we view $Y_0$ and $Y_\infty$ as closed subschemes of $X$
via the morphism $p$ then we have
$$
\text{div}_X(f) = [Y_0]_{n - 1} - [Y_\infty]_{n - 1}
$$
\end{enumerate}
\end{lemma}

\begin{proof}
Since $X$ is integral, we see that $U$ is integral.
Hence $Y$ is integral, and $(1, f)(U) \subset Y$ is an open dense subscheme.
Also, note that the closed subscheme $Y \subset X \times_S \mathbf{P}^1_S$
does not depend on the choice of the open $U$, since after all it is
the closure of the one point set $\{\eta'\} = \{(1, f)(\eta)\}$
where $\eta \in X$ is the generic point. Having said this let us
prove the assertions of the lemma.

\medskip\noindent
For (1) note that $p$ is the composition of the closed immersion
$Y \to X \times_S \mathbf{P}^1_S = \mathbf{P}^1_X$ with the proper
morphism $\mathbf{P}^1_X \to X$. As a composition of proper morphisms
is proper (Morphisms, Lemma \ref{morphisms-lemma-composition-proper})
we conclude.

\medskip\noindent
It is clear that $Y \cap U \times_S \mathbf{P}^1_S = (1, f)(U)$.
Thus (2) follows. It also follows that $\dim_\delta(Y) = n$.

\medskip\noindent
Note that $q(\eta') = f(\eta)$ is not contained in $D_0$ or $D_\infty$
since $f \in R(X)^*$. Hence (3) by
Divisors, Lemma \ref{divisors-lemma-pullback-effective-Cartier-defined}.
We obtain $\dim_\delta(Y_0) = n - 1$
and $\dim_\delta(Y_\infty) = n - 1$ from
Lemma \ref{lemma-divisor-delta-dimension}.

\medskip\noindent
Consider the effective Cartier divisor $Y_0$.
At every point $\xi \in Y_0$ we have $f \in \mathcal{O}_{Y, \xi}$ and
the local equation for $Y_0$ is given by $f$.
In particular, if $\delta(\xi) = n - 1$ so $\xi$ is the generic point
of an integral closed subscheme $Z$ of $\delta$-dimension $n - 1$,
then we see that the coefficient of $[Z]$ in $\text{div}_Y(f)$ is
$$
\text{ord}_Z(f) =
\text{length}_{\mathcal{O}_{Y, \xi}}
(\mathcal{O}_{Y, \xi}/f\mathcal{O}_{Y, \xi}) =
\text{length}_{\mathcal{O}_{Y, \xi}}
(\mathcal{O}_{Y_0, \xi})
$$
which is the coefficient of $[Z]$ in $[Y_0]_{n - 1}$. A similar
argument using the rational function $1/f$ shows that
$-[Y_\infty]$ agrees with the terms with negative coefficients in
the expression for $\text{div}_Y(f)$. Hence (4) follows.

\medskip\noindent
Note that $D_0 \to S$ is an isomorphism. Hence we see that
$X \times_S D_0 \to X$ is an isomorphism as well. Clearly
we have $Y_0 = Y \cap X \times_S D_0$ (scheme theoretic intersection)
inside $X \times_S \mathbf{P}^1_S$. Hence it is really the case that
$Y_0 \to X$ is a closed immersion. It follows that
$$
p_*\mathcal{O}_{Y_0} = \mathcal{O}_{Y'_0}
$$
where $Y'_0 \subset X$ is the image of $Y_0 \to X$.
By Lemma \ref{lemma-cycle-push-sheaf} we
have $p_*[Y_0]_{n - 1} = [Y'_0]_{n - 1}$. The same
is true for $D_\infty$ and $Y_\infty$. Hence (6) is a consequence of (5).
Finally, (5) follows immediately from
Lemma \ref{lemma-proper-pushforward-alteration}.
\end{proof}

\end{document}
